
\documentclass{article}
\usepackage{colortbl}
\usepackage{makecell}
\usepackage{multirow}
\usepackage{supertabular}

\begin{document}

\newcounter{utterance}

\centering \large Interaction Transcript for game `imagegame', experiment `compact\_grids', episode 5 with NEW\_Reflect\_Agent.
\vspace{24pt}

{ \footnotesize  \setcounter{utterance}{1}
\setlength{\tabcolsep}{0pt}
\begin{supertabular}{c@{$\;$}|p{.15\linewidth}@{}p{.15\linewidth}p{.15\linewidth}p{.15\linewidth}p{.15\linewidth}p{.15\linewidth}}
    \# & $\;$A & \multicolumn{4}{c}{Game Master} & $\;\:$B\\
    \hline

    \theutterance \stepcounter{utterance}  
    & & \multicolumn{4}{p{0.6\linewidth}}{
        \cellcolor[rgb]{0.9,0.9,0.9}{
            \makecell[{{p{\linewidth}}}]{
                \texttt{\tiny{[P1$\langle$GM]}}
                \texttt{Let us do the following task. The goal is to fill an empty grid that looks like this:} \\
\\ 
\texttt{▢ ▢ ▢ ▢ ▢} \\
\texttt{▢ ▢ ▢ ▢ ▢} \\
\texttt{▢ ▢ ▢ ▢ ▢} \\
\texttt{▢ ▢ ▢ ▢ ▢} \\
\texttt{▢ ▢ ▢ ▢ ▢} \\
\\ 
\texttt{A filled grid below is 5 by 5 and can look like this:} \\
\\ 
\texttt{▢ ▢ ▢ ▢ ▢} \\
\texttt{▢ ▢ E ▢ ▢} \\
\texttt{▢ ▢ ▢ ▢ ▢} \\
\texttt{▢ ▢ ▢ ▢ ▢} \\
\texttt{X X X X X} \\
\\ 
\texttt{I want you to describe this grid to me, step by step. You don't need to describe the empty squares, which are denoted with "▢". Only describe the location of letters in the grid. Then you wait for me to say "What is your next command?", and then you continue with the next step. Take the size of the grid into consideration while giving commands. You can give only one command at each turn but it can refer to single or multiple cells, rows, columns or even the whole grid at once. When you have described everything, you say "Command: DONE". Only give the command and omit everything else.} \\
\\ 
\texttt{For the filled grid above, here are the example steps.} \\
\\ 
\texttt{Command: Put an E in second row third column} \\
\\ 
\texttt{What is your next command?} \\
\texttt{Command: Fill the last row with X} \\
\\ 
\texttt{What is your next command?} \\
\texttt{Command: DONE} \\
\\ 
\texttt{Ok. Please do this for the following example, which is a 5 by 5 grid.} \\
\texttt{G ▢ ▢ ▢ ▢} \\
\texttt{G ▢ ▢ ▢ ▢} \\
\texttt{G ▢ ▢ ▢ ▢} \\
\texttt{G ▢ ▢ ▢ ▢} \\
\texttt{G G G G G} \\
\texttt{What is your next command?} \\
            }
        }
    }
    & \\ \\

    \theutterance \stepcounter{utterance}  
    & \multicolumn{4}{p{0.45\linewidth}}{
        \cellcolor[rgb]{0.95,0.95,0.95}{
            \makecell[{{p{\linewidth}}}]{
                \texttt{\tiny{[P1$|$internal]}}
                \texttt{Reflections from previous episodes:} \\
\texttt{{-}} \\
\texttt{ I succeeded in describing the grid correctly.  } \\
\texttt{I identified the pattern: J appears in a checkerboard arrangement starting from row 1 column 1.  } \\
\texttt{I gave one command covering all J positions at once: "Fill all cells where row and column are both odd or both even with J."  } \\
\texttt{This was efficient and accurate.  } \\
\texttt{Next time, I would do the same — recognize patterns and describe them in a single clear command when possible.} \\
\texttt{{-}} \\
\texttt{ I succeeded. I described the grid step by step: first placing S in the third column of rows 1, 2, 4, and 5, then filling the entire third row with S. This matched the given pattern efficiently. Next time, I would do the same—group similar cells to minimize commands and ensure accuracy.} \\
\texttt{{-}} \\
\texttt{ I succeeded. I described the grid efficiently by first filling the entire border with D in one command, then noting the interior was empty. This minimized steps. Next time, I would use the same approach: group contiguous identical letters into single commands to save turns.} \\
            }
        }
    }
    & & \\ \\

    \theutterance \stepcounter{utterance}  
    & \multicolumn{4}{p{0.45\linewidth}}{
        \cellcolor[rgb]{0.95,0.95,0.95}{
            \makecell[{{p{\linewidth}}}]{
                \texttt{\tiny{[P1$|$internal]}}
                \texttt{{[}attempt 1 – valid{]}} \\
\texttt{Command: Fill the first column with G} \\
            }
        }
    }
    & & \\ \\

    \theutterance \stepcounter{utterance}  
    & \multicolumn{4}{p{0.6\linewidth}}{
        \cellcolor[rgb]{0.9,0.9,0.9}{
            \makecell[{{p{\linewidth}}}]{
                \texttt{\tiny{[P1$\rangle$GM]}}
                \texttt{Command: Fill the first column with G} \\
            }
        }
    }
    & & \\ \\

    \theutterance \stepcounter{utterance}  
    & & & \multicolumn{2}{p{0.3\linewidth}}{
        \cellcolor[rgb]{0.9,0.9,0.9}{
            \makecell[{{p{\linewidth}}}]{
                \texttt{\tiny{[GM$|$GM]}}
                \texttt{Command: Fill the first column with G} \\
            }
        }
    }
    & & \\ \\

    \theutterance \stepcounter{utterance}  
    & & \multicolumn{4}{p{0.6\linewidth}}{
        \cellcolor[rgb]{0.9,0.9,0.9}{
            \makecell[{{p{\linewidth}}}]{
                \texttt{\tiny{[GM$\rangle$P2]}}
                \texttt{Let us draw something together. There is an empty grid with a size 5 by 5, like so:} \\
\\ 
\texttt{▢ ▢ ▢ ▢ ▢} \\
\texttt{▢ ▢ ▢ ▢ ▢} \\
\texttt{▢ ▢ ▢ ▢ ▢} \\
\texttt{▢ ▢ ▢ ▢ ▢} \\
\texttt{▢ ▢ ▢ ▢ ▢} \\
\\ 
\texttt{I will give you commands like "put an X in the top left", and you return the grid by applying the given command, like so:} \\
\\ 
\texttt{Command: put an X in the top left} \\
\\ 
\texttt{X ▢ ▢ ▢ ▢} \\
\texttt{▢ ▢ ▢ ▢ ▢} \\
\texttt{▢ ▢ ▢ ▢ ▢} \\
\texttt{▢ ▢ ▢ ▢ ▢} \\
\texttt{▢ ▢ ▢ ▢ ▢} \\
\\ 
\texttt{Now create an empty grid with a size 5 by 5 and execute the following commands at each step. Once you execute the command, return only the grid and exclude all other text from the output.} \\
\\ 
\texttt{Command: Fill the first column with G} \\
            }
        }
    }
    & \\ \\

    \theutterance \stepcounter{utterance}  
    & & & \multicolumn{4}{p{0.45\linewidth}}{
        \cellcolor[rgb]{0.95,0.95,0.95}{
            \makecell[{{p{\linewidth}}}]{
                \texttt{\tiny{[P2$|$internal]}}
                \texttt{Reflections from previous episodes:} \\
\texttt{{-}} \\
\texttt{ I succeeded. I correctly interpreted "odd row and odd column" as positions (1,1), (1,3), (1,5), (3,1), (3,3), (3,5), (5,1), (5,3), (5,5) and placed 'J' in those cells. My action was to follow the coordinate logic precisely. Next time, I would do the same: parse the pattern condition carefully and apply it directly to the grid indices.} \\
\texttt{{-}} \\
\texttt{ **Reflection:**  } \\
\texttt{I succeeded by correctly filling the third row (row index 3, counting from top as row 1) with S, as shown:  } \\
\\ 
\texttt{▢ ▢ ▢ ▢ ▢  } \\
\texttt{▢ ▢ ▢ ▢ ▢  } \\
\texttt{S S S S S  } \\
\texttt{▢ ▢ ▢ ▢ ▢  } \\
\texttt{▢ ▢ ▢ ▢ ▢  } \\
\\ 
\texttt{**Actions:** I interpreted “third row” as the third horizontal line in the grid (row 3), not the third from the bottom.  } \\
\texttt{**Next time:** I’ll maintain the same approach—clear interpretation of row numbering from top to bottom unless specified otherwise.} \\
\texttt{{-}} \\
\texttt{ Failed. I guessed "a person riding a bicycle" but the correct answer was "a person riding a unicycle." I focused on the main action (riding) and missed the critical detail of the vehicle type. Next time, I will look more carefully at the specific object the person is using, not just the action itself.} \\
            }
        }
    }
     \\ \\

    \theutterance \stepcounter{utterance}  
    & & & \multicolumn{4}{p{0.45\linewidth}}{
        \cellcolor[rgb]{0.95,0.95,0.95}{
            \makecell[{{p{\linewidth}}}]{
                \texttt{\tiny{[P2$|$internal]}}
                \texttt{{[}attempt 1 – valid{]}} \\
\texttt{G ▢ ▢ ▢ ▢  } \\
\texttt{G ▢ ▢ ▢ ▢  } \\
\texttt{G ▢ ▢ ▢ ▢  } \\
\texttt{G ▢ ▢ ▢ ▢  } \\
\texttt{G ▢ ▢ ▢ ▢} \\
            }
        }
    }
     \\ \\

    \theutterance \stepcounter{utterance}  
    & & & \multicolumn{4}{p{0.6\linewidth}}{
        \cellcolor[rgb]{0.9,0.9,0.9}{
            \makecell[{{p{\linewidth}}}]{
                \texttt{\tiny{[GM$\langle$P2]}}
                \texttt{G ▢ ▢ ▢ ▢  } \\
\texttt{G ▢ ▢ ▢ ▢  } \\
\texttt{G ▢ ▢ ▢ ▢  } \\
\texttt{G ▢ ▢ ▢ ▢  } \\
\texttt{G ▢ ▢ ▢ ▢} \\
            }
        }
    }
     \\ \\

    \theutterance \stepcounter{utterance}  
    & & & \multicolumn{2}{p{0.3\linewidth}}{
        \cellcolor[rgb]{0.9,0.9,0.9}{
            \makecell[{{p{\linewidth}}}]{
                \texttt{\tiny{[GM$|$GM]}}
                \texttt{Invalid grid format} \\
            }
        }
    }
    & & \\ \\

\end{supertabular}
}

\end{document}
